\section{Evaluation}
\label{sec:eval}

\subsection{Experimental Setup}
\label{subsec:setup}

% VALUE: one setup paragraph: platform, software and versions, workloads, the baseline, how TTFT is measured, the correctness gate, and how the machine state is controlled; defaults stated once.
\textbf{Platform.}
We measure on a GMKtec NucBox EVO-X2 with an AMD Ryzen\texttrademark{} AI Max+ 395 processor, whose XDNA~2 NPU has 8 columns of 4 compute tiles, and 64\,GB of shared system memory.
The host runs Linux~7.0 with the amdxdna driver and XRT~2.25.0, the NPU in its performance power mode, and the CPU's frequency governor and energy preference set to performance.
% uname -r 7.0.0-34-generic; amdxdna 2.25.0_20260922 (DKMS, 31c01f2019e7); XRT 2.25.0 (3923b2b5be55); BIOS EVO-X2 1.12; NPU firmware version: \E{record}; read 2026-10-14. No OS platform-profile interface exists on this machine (lab-prefill-overlap (11)); the BIOS power mode is \E{record}.

\textbf{Workloads.}
We evaluate prefill of Llama-3.2-1B~\cite{meta2024llama32} and Qwen3-4B~\cite{yang2025qwen3} with 4-bit quantized weights in FastFlowLM's format~\cite{fastflowlm-repo}, at prompts of 512, 2{,}048, and 4{,}096 tokens.
\E{Further models and the decode and agent-step workloads, once measured.}
The designs are MLIR-AIR's fused prefill~\cite{wang2026mlirair}, compiled with the MLIR-AIE compiler~\cite{mlir-aie-repo}, and are used unchanged.
% mlir-air a3e26261 (llms/shared/fused_prefill); mlir_aie wheel 1.4.5.dev22+g35fa6b1. Weights: FastFlowLM/Qwen3-4B-NPU2 rev 07651f11cac7 (model.q4nx sha256-verified); Llama-3.2-1B q4nx sha256 80a05efc... (lab-prefill-overlap (23)).

\textbf{Baseline and Correctness.}
The baseline is each design's own host program, which runs the same configuration, kernels, and weights as our composed schedules.
Every run's logits must be bit-identical to the baseline's, and the kernel log must show no driver or IOMMU message during the run.

\textbf{Measurement.}
TTFT is the wall time of the prefill call, from token identifiers to logits, and we report the median of 3 timed runs after one warm run, with the order of schedules rotated between rounds.
Each batch of runs holds an exclusive lock on the device, checks that no other process holds it, and brackets the runs with a DRAM bandwidth probe and an empty-command dispatch measurement.
The aggregate DRAM read cap varies between 50 and 56\,GB/s across batches, so every ratio compares runs within one batch.
% Lock: flock ~/.local/state/npu.lock; brackets: roofline.py 1x1 / 8x2 reads (256 MiB), dispatch floor (4 KB bypass), loadavg, NPU mode, EPP, governor (lab-prefill-device (1), lab-prefill-overlap (2), (11)).

% ---------------------------------------------------------------------------------------------------------------
% Results subsections: placeholders until the decode-half and agent-step results land and the headline numbers pass
% the independent audit (hub Status). Each subsection opens with the pattern it supports and ends with it.

\subsection{Prefill Latency}
\label{subsec:ttft}
% VALUE: composing concurrent commands lowers TTFT on every model and length, more at longer prompts where attention's share grows; report the essential decisions (co-scheduled groups, instruction order) separately from the host-boundary handling (queue depth, packing).
\E{TTFT by model and prompt length; the ladder from the existing schedule through co-scheduled groups, completion order, and host-boundary handling.}

\subsection{Generality Across Models}
\label{subsec:generality}
% VALUE: the same composer and policy rule apply unchanged to a second model family, with bit-identical outputs.
\E{Second and further models.}

\subsection{Model Accuracy}
\label{subsec:accuracy}
% VALUE: the event model predicts held-out schedules and an unseen model within a few percent, component by component, and states where it misses (longer prompts) and why.
\E{Held-out and unseen-model predictions, per component; the regime where the model misses.}

\subsection{Decode and Agent Steps}
\label{subsec:decode}
% VALUE: pending the decode-half and agent-step measurements.
\E{Decode run-ahead and end-to-end agent-shaped turns.}
